\documentclass[12pt,a4paper]{article}
\usepackage[utf8]{inputenc}
\usepackage[german]{babel}
\usepackage{amsmath}
\usepackage{amssymb}
\usepackage{graphicx}
\usepackage{tikz}
\usepackage{xcolor}
\usepackage{fancyhdr}
\usepackage{geometry}
\usetikzlibrary{positioning}
\geometry{margin=2cm}
\setlength{\headheight}{14.49998pt}

\pagestyle{fancy}
\fancyhf{}
\rhead{Klausur Datenbanken I - Musterlösung}
\lhead{DHBW TINF19E}
\cfoot{\thepage}

\title{\textbf{Klausur Datenbanken I - Musterlösung}}
\author{Modul T2INF2004 \\[0.3cm] Dozent: Stark \\[0.3cm] Datum: 25.05.2021}
\date{Bearbeitungszeit: 60 Minuten}

\begin{document}

\maketitle
\vspace{1cm}

\noindent
\textbf{Gesamtpunkte:} 42 Punkte

\vspace{1.5cm}

\section*{Aufgabenteil 1 - Grundlagen (13 Punkte)}

\subsection*{1. Hierarchisches vs. Netzwerkmodell}

\textbf{Problem der hierarchischen Datenbank:}
Das hierarchische Modell konnte nur 1:n-Beziehungen abbilden, nicht aber m:n-Beziehungen. Dies führte zu Redundanzen und Anomalien.

\textbf{Lösung des Netzwerkmodells:}
Das Netzwerkmodell ermöglichte m:n-Beziehungen und beseitigte damit die Redundanzen des hierarchischen Modells.

\textbf{Warum hat es sich nicht durchgesetzt:}
\begin{itemize}
    \item Hohe Komplexität in der Anwendung und Navigation
    \item Manuelle Navigation durch Zeiger und Verweise notwendig
    \item Mangelnde Benutzerfreundlichkeit
    \item Das relationale Modell bot eine einfachere, intuitivere Alternative mit SQL
\end{itemize}

\subsection*{2. Sequentielle Durchläufe im Datenbankentwurfsprozess}

\textbf{Warum sequenziell:}
\begin{itemize}
    \item Jede Stufe baut auf den Ergebnissen der vorherigen auf
    \item Festlegung von Anforderungen ist Basis für konzeptionellen Entwurf
    \item Logisches Modell basiert auf konzeptionellem Modell
    \item Physischer Entwurf erfordert logisches Modell
\end{itemize}

\textbf{Problem bei Auslassung des konzeptionellen Entwurfs:}
\begin{itemize}
    \item Keine klare Modellierung der realen Welt
    \item Fehlende Semantic führt zu inkonsistentem relationalem Modell
    \item Redundanzen und Anomalien nicht vermieden
    \item Spätere Änderungen führen zu großen Umstrukturierungen
    \item Keine dokumentierte Grundlage für Wartung und Evolution
\end{itemize}

\subsection*{3. Künstliche Schlüssel im ERM}

\textbf{Wann sinnvoll:}
\begin{itemize}
    \item Wenn kein natürlicher Schlüssel existiert oder dieser zu komplex wäre
    \item Bei Entitäten mit mehreren möglichen Schlüsseln (Designentscheidung)
    \item Zur Vereinfachung und Leistungsoptimierung
\end{itemize}

\textbf{Realisierung in der DDL:}
\begin{verbatim}
CREATE TABLE Kunde (
    Kunden_ID INT PRIMARY KEY AUTO_INCREMENT,
    Name VARCHAR(100),
    Adresse VARCHAR(200)
);
\end{verbatim}

Der künstliche Schlüssel wird als \texttt{PRIMARY KEY} mit \texttt{AUTO\_INCREMENT} definiert.

\subsection*{4. Existenzabhängige Entität in DDL}

\textbf{1. Variante: Fremdschlüssel mit NOT NULL und ON DELETE CASCADE}
\begin{verbatim}
CREATE TABLE Bestellung (
    Bestellung_ID INT,
    Kunde_ID INT NOT NULL,
    Datum DATE,
    PRIMARY KEY (Bestellung_ID, Kunde_ID),
    FOREIGN KEY (Kunde_ID) REFERENCES Kunde(Kunde_ID)
        ON DELETE CASCADE
);
\end{verbatim}

\textbf{2. Variante: Fremdschlüssel als Teil des Primärschlüssels}
\begin{verbatim}
CREATE TABLE Bestellposition (
    Bestellung_ID INT NOT NULL,
    Position_Nr INT NOT NULL,
    Artikel_ID INT,
    Menge INT,
    PRIMARY KEY (Bestellung_ID, Position_Nr),
    FOREIGN KEY (Bestellung_ID) REFERENCES Bestellung(Bestellung_ID)
);
\end{verbatim}

Bei Variante 1 wird die Existenzabhängigkeit durch NOT NULL erzwungen. Bei Variante 2 ist die abhängige Entität nicht selbstständig identifizierbar.

\subsection*{5. Prim-Attribut in der 3NF-Prüfung}

\textbf{Definition:}
Ein Prim-Attribut ist ein Attribut, das Teil eines Schlüsselkandidaten ist.

\textbf{Verwendung bei 3NF-Prüfung:}
\begin{itemize}
    \item 3NF verlangt: Jedes Nicht-Schlüssel-Attribut (Non-Prime-Attribut) darf nicht transitiv vom Schlüssel abhängen
    \item Prim-Attribute sind von dieser Regel ausgenommen
    \item Die Prüfung konzentriert sich auf Abhängigkeiten von Non-Prime-Attributen
    \item Ziel: Vermeidung von transitiven Abhängigkeiten (A → B → C)
\end{itemize}

\subsection*{6. PostgreSQL Query - Elemente ohne Beziehung}

\textbf{Problem:} Alle Elemente aus Tabelle A anzeigen, die \textbf{nicht} an einer Beziehung mit B beteiligt sind.

\textbf{Lösung:}
\begin{verbatim}
SELECT * FROM A
WHERE A.id NOT IN (SELECT id FROM B);
\end{verbatim}

Oder mit LEFT OUTER JOIN:
\begin{verbatim}
SELECT A.* FROM A
LEFT OUTER JOIN B ON A.id = B.id
WHERE B.id IS NULL;
\end{verbatim}

Die zweite Variante ist in der Regel effizienter.

\vspace{2cm}

\section*{Aufgabenteil 2 - ER-Modell (12 Punkte)}

\subsection*{a) ER-Modell}

\subsubsection*{Szenario-Analyse und Anforderungen}

Das Szenario beschreibt ein Fahrzeugvermietungssystem für Reiseveranstalter. Folgende Anforderungen lassen sich extrahieren:

\begin{enumerate}
    \item \textbf{Buchungsprozess:} Ein Kunde bucht Mietfahrzeuge für eine Wegstrecke (Hin- und Rückfahrt)
    \item \textbf{Fahrzeugauswahl:} Für eine Wegstrecke können mehrere verschiedene Fahrzeugtypen gebucht werden
    \item \textbf{Mehrfachbuchungen:} Mehrere Fahrzeuge können für eine Strecke eingesetzt werden
    \item \textbf{Anbieter-Beschränkung:} Pro Wegstrecke wird aber nur ein Anbieter verwendet (eine Buchung = ein Fahrzeug bei einem Anbieter)
    \item \textbf{Tarifverwaltung:} Jeder Anbieter definiert seine eigenen Tarife
    \item \textbf{Rechnungsverwaltung:} Mehrere Buchungen einer Wegstrecke werden auf einer Rechnung zusammengefasst
\end{enumerate}

\subsubsection*{Identifizierte Entitäten und deren Begründung}

\begin{enumerate}
    \item \textbf{Kunde}
    \begin{itemize}
        \item Attributes: Kunde\_ID (PK), Name, Adresse, Telefon
        \item Begründung: Zentrale Entität; bucht Fahrzeuge
    \end{itemize}

    \item \textbf{Wegstrecke}
    \begin{itemize}
        \item Attributes: Strecken\_ID (PK), Startort, Zielort, Entfernung, Datum\_Start, Datum\_Ende
        \item Begründung: Repräsentiert eine Hin- und Rückfahrtstrecke; mehrere Fahrzeuge möglich
    \end{itemize}

    \item \textbf{Anbieter}
    \begin{itemize}
        \item Attributes: Anbieter\_ID (PK), Name, Adresse, Telefon, Email
        \item Begründung: Vermietungsunternehmen; hat Fahrzeuge und Tarife
    \end{itemize}

    \item \textbf{Fahrzeug} (existenzabhängig von Anbieter)
    \begin{itemize}
        \item Attributes: Fahrzeug\_ID, Kennzeichen, Typ, Baujahr, Anbieter\_ID (FK)
        \item Composite Key: (Fahrzeug\_ID, Anbieter\_ID)
        \item Begründung: Gehört immer zu einem Anbieter; ohne Anbieter keine Existenz möglich
    \end{itemize}

    \item \textbf{Tarif} (existenzabhängig von Anbieter)
    \begin{itemize}
        \item Attributes: Tarif\_ID, Bezeichnung, Preis\_pro\_Tag, Kategorien, Anbieter\_ID (FK)
        \item Composite Key: (Tarif\_ID, Anbieter\_ID)
        \item Begründung: Tarife sind anbieter-spezifisch; ohne Anbieter nicht sinnvoll
    \end{itemize}

    \item \textbf{Buchung}
    \begin{itemize}
        \item Attributes: Buchungs\_ID (PK), Kunde\_ID (FK), Wegstrecke\_ID (FK), Fahrzeug\_ID (FK), Anbieter\_ID (FK), Tarif\_ID (FK), Buchungsdatum, Mietdauer, Kosten
        \item Begründung: Assoziative Entität; verbindet Kunde, Strecke, Fahrzeug und Tarif; enthält Buchungsdetails
        \item Multiplicities: Ein Kunde viele Buchungen, eine Strecke viele Buchungen, ein Fahrzeug viele Buchungen
    \end{itemize}

    \item \textbf{Rechnung} (existenzabhängig von Buchung - ODER direkt von Strecke)
    \begin{itemize}
        \item Attributes: Rechnungs\_ID (PK), Strecken\_ID (FK), Kunde\_ID (FK), Rechnungsdatum, Gesamtbetrag, Status
        \item Begründung: Fasst mehrere Buchungen einer Strecke zusammen; existiert nur, wenn Buchungen vorhanden
    \end{itemize}
\end{enumerate}

\subsubsection*{ER-Diagramm}

\begin{figure}[h]
\centering
\begin{tikzpicture}[scale=0.75, node distance=2cm]
  % Styles
  \tikzset{
    entity/.style={rectangle, draw=black, very thick, minimum width=2cm, minimum height=1cm},
    weak entity/.style={entity, double, double distance=2pt},
    relation/.style={diamond, draw=black, very thick, minimum width=1.5cm, minimum height=1cm}
  }
  
  % Entities
  \node[entity] (Kunde) at (0, 6) {KUNDE};
  \node[below=0.1cm of Kunde.south, inner sep=0pt] {\tiny Kunde\_ID, Name};
  
  \node[entity] (Strecke) at (3, 6) {WEGSTRECKE};
  \node[below=0.1cm of Strecke.south, inner sep=0pt] {\tiny Strecken\_ID, Ort};
  
  \node[entity] (Anbieter) at (6, 6) {ANBIETER};
  \node[below=0.1cm of Anbieter.south, inner sep=0pt] {\tiny Anbieter\_ID, Name};
  
  \node[entity] (Buchung) at (3, 3) {BUCHUNG};
  \node[below=0.1cm of Buchung.south, inner sep=0pt] {\tiny Buchungs\_ID, Kosten};
  
  \node[weak entity] (Fahrzeug) at (6, 3) {FAHRZEUG};
  \node[below=0.1cm of Fahrzeug.south, inner sep=0pt] {\tiny Fahrzeug\_ID, Typ};
  
  \node[weak entity] (Tarif) at (7.5, 4.5) {TARIF};
  \node[below=0.1cm of Tarif.south, inner sep=0pt] {\tiny Tarif\_ID, Preis};
  
  \node[weak entity] (Rechnung) at (1.5, 1) {RECHNUNG};
  \node[below=0.1cm of Rechnung.south, inner sep=0pt] {\tiny Rechnungs\_ID};
  
  % Relationships
  \draw (Kunde) -- node[above, font=\small] {bucht} (Buchung);
  \draw (Strecke) -- node[above left, font=\small] {fuer} (Buchung);
  \draw (Fahrzeug) -- node[right, font=\small] {mit} (Buchung);
  \draw (Tarif) -- node[right, font=\small] {nach} (Buchung);
  \draw (Anbieter) -- node[above, font=\small] {besitzt} (Fahrzeug);
  \draw (Anbieter) -- node[left, font=\small] {hat} (Tarif);
  \draw (Buchung) -- node[left, font=\small] {enthaelt} (Rechnung);
  
  % Cardinality labels
  \node[font=\tiny] at (0.8, 5) {1:n};
  \node[font=\tiny] at (2, 5) {1:n};
  \node[font=\tiny] at (5, 3.7) {1:n};
  \node[font=\tiny] at (6.5, 4) {1:n};
  \node[font=\tiny] at (6.2, 5) {1:n};
  \node[font=\tiny] at (7, 4) {1:n};
  \node[font=\tiny] at (2, 2) {1:1};
\end{tikzpicture}
\caption{ER-Diagramm des Fahrzeugvermietungssystems}
\end{figure}

\subsubsection*{Detaillierte Beziehungsbeschreibung}

\begin{enumerate}
    \item \textbf{Kunde -- Buchung (1:n)}
    \begin{itemize}
        \item Ein Kunde kann viele Buchungen tätigen
        \item Eine Buchung gehört genau zu einem Kunde
    \end{itemize}

    \item \textbf{Wegstrecke -- Buchung (1:n)}
    \begin{itemize}
        \item Eine Wegstrecke kann mehrere Buchungen haben (verschiedene Fahrzeuge)
        \item Eine Buchung bezieht sich auf genau eine Wegstrecke
    \end{itemize}

    \item \textbf{Fahrzeug -- Buchung (1:n) [mit Integrität]}
    \begin{itemize}
        \item Ein Fahrzeug kann für mehrere Strecken gebucht werden
        \item Wichtig: Eine Buchung = ein Fahrzeug bei einem Anbieter (Constraint)
    \end{itemize}

    \item \textbf{Tarif -- Buchung (1:n)}
    \begin{itemize}
        \item Ein Tarif kann für mehrere Buchungen gelten
        \item Eine Buchung nutzt genau einen Tarif
    \end{itemize}

    \item \textbf{Anbieter -- Fahrzeug (1:n) [Existenzabhängigkeit]}
    \begin{itemize}
        \item Ein Anbieter besitzt viele Fahrzeuge
        \item Ein Fahrzeug gehört genau zu einem Anbieter
        \item Existenzabhängigkeit: Fahrzeug kann ohne Anbieter nicht existieren
    \end{itemize}

    \item \textbf{Anbieter -- Tarif (1:n) [Existenzabhängigkeit]}
    \begin{itemize}
        \item Ein Anbieter hat mehrere Tarife
        \item Ein Tarif gehört genau zu einem Anbieter
        \item Existenzabhängigkeit: Tarif ist anbieterspezifisch
    \end{itemize}

    \item \textbf{Buchung -- Rechnung (1:1 oder 1:n Abhängigkeit)}
    \begin{itemize}
        \item Mehrere Buchungen einer Strecke werden auf einer Rechnung zusammengefasst
        \item Modellierungsvariante: Rechnung ist existenzabhängig (doppeltes Rechteck)
        \item Bessere Variante: Rechnung direkt mit Strecke verbinden (1:1)
    \end{itemize}
\end{enumerate}

\subsection*{b) Löschregeln für existenzabhängige Entitäten}

\subsubsection*{Allgemeines Konzept der Löschregeln}

Existenzabhängige Entitäten werden durch Löschregeln charakterisiert:
\begin{itemize}
    \item \textbf{Cascade (C):} Wenn die übergeordnete Entität gelöscht wird, werden auch alle abhängigen Entitäten gelöscht
    \item \textbf{Restrict (R):} Das Löschen ist nicht erlaubt, solange abhängige Entitäten existieren
    \item \textbf{Set Null (N):} Der Fremdschlüssel wird auf NULL gesetzt (nicht bei Existenzabhängigkeit!)
\end{itemize}

\subsubsection*{Konkrete Löschregeln für das Szenario}

\textbf{1. Fahrzeug (abhängig von Anbieter) - Regel: RESTRICT (R)}

\textbf{Begründung:}
\begin{itemize}
    \item Ein Fahrzeug gehört immer zu genau einem Anbieter
    \item Fahrzeuge sind oft in laufende Buchungen involviert
    \item Ein automatisches Löschen würde zu Integritätsverletzungen führen
    \item Empfehlung: RESTRICT, da Fahrzeuge archiviert statt gelöscht werden sollten
\end{itemize}

\textbf{DDL-Implementierung:}
\begin{verbatim}
CREATE TABLE Fahrzeug (
    Fahrzeug_ID INT NOT NULL,
    Anbieter_ID INT NOT NULL,
    Kennzeichen VARCHAR(10),
    Typ VARCHAR(50),
    PRIMARY KEY (Fahrzeug_ID, Anbieter_ID),
    FOREIGN KEY (Anbieter_ID) 
        REFERENCES Anbieter(Anbieter_ID)
        ON DELETE RESTRICT
);
\end{verbatim}

\textbf{2. Tarif (abhängig von Anbieter) - Regel: RESTRICT (R)}

\textbf{Begründung:}
\begin{itemize}
    \item Tarife sind essentiell für die Kostenberechnung von Buchungen
    \item Historische Daten müssen erhalten bleiben (Audit Trail)
    \item Ein Löschen würde Rechnungen ungültig machen
    \item Empfehlung: RESTRICT oder ein Flag ''aktiv/inaktiv'' statt Löschen
\end{itemize}

\textbf{DDL-Implementierung:}
\begin{verbatim}
CREATE TABLE Tarif (
    Tarif_ID INT NOT NULL,
    Anbieter_ID INT NOT NULL,
    Bezeichnung VARCHAR(100),
    Preis_pro_Tag DECIMAL(10, 2),
    PRIMARY KEY (Tarif_ID, Anbieter_ID),
    FOREIGN KEY (Anbieter_ID) 
        REFERENCES Anbieter(Anbieter_ID)
        ON DELETE RESTRICT
);
\end{verbatim}

\textbf{3. Rechnung (abhängig von Buchung/Strecke) - Regel: RESTRICT (R) oder CASCADE (C)}

\textbf{Möglichkeit A: RESTRICT (R) - Konservativ}
\begin{itemize}
    \item Eine Rechnung muss erhalten bleiben (Buchhaltung!)
    \item Löschen ist nicht sinnvoll
    \item Empfehlung: RESTRICT + Status-Flag (''storniert'')
\end{itemize}

\textbf{Möglichkeit B: CASCADE (C) - Mit Vorsicht}
\begin{itemize}
    \item Falls Buchung storniert wird, kann die Rechnung mitgelöscht werden
    \item Aber: Rechnungen sollten aus Compliance-Gründen erhalten bleiben!
    \item Besser: Soft-Delete mit Status-Feld
\end{itemize}

\textbf{Empfehlung:} CASCADE ist hier nicht sinnvoll. Besser: RESTRICT + Stornierungslogik.

\textbf{DDL-Implementierung (mit Statusfeld):}
\begin{verbatim}
CREATE TABLE Rechnung (
    Rechnungs_ID INT PRIMARY KEY AUTO_INCREMENT,
    Strecken_ID INT NOT NULL,
    Kunde_ID INT NOT NULL,
    Gesamtbetrag DECIMAL(10, 2),
    Status VARCHAR(20) DEFAULT 'offen',
    -- Status: 'offen', 'bezahlt', 'storniert'
    FOREIGN KEY (Strecken_ID) 
        REFERENCES Wegstrecke(Strecken_ID),
    FOREIGN KEY (Kunde_ID) 
        REFERENCES Kunde(Kunde_ID)
);
\end{verbatim}

\vspace{2cm}

\section*{Aufgabenteil 3 - Relationale Algebra und SQL (8 Punkte)}

\subsection*{a) Relationale Algebra}

\textbf{Aufgabe:} Welche Ärzte arbeiten auf Station 3?

\textbf{Lösung:}
\begin{align*}
\text{Ergebnis} = & \pi_{\text{Pers-Nr, Nachname, Vorname, Fachgebiet, Rang}} \\
& \left( \sigma_{\text{Stat-Nr} = 3} (\text{Zimmer}) \right) \\
& \bowtie (\text{Patient} \bowtie \text{behandelt} \bowtie \text{Ärzte} \bowtie \text{Stations\_Personal})
\end{align*}

\textbf{Alternative ohne Join-Operator (nur Basis-Operatoren):}
\begin{align*}
\text{Step 1:} & \quad Z_3 = \sigma_{\text{Stat-Nr} = 3} (\text{Zimmer}) \\
\text{Step 2:} & \quad P_3 = \sigma_{\text{Raum-Nr} \in Z_3.\text{Raum-Nr}} (\text{Patient}) \\
\text{Step 3:} & \quad B_3 = \sigma_{\text{Pat-Nr} \in P_3.\text{Pat-Nr}} (\text{behandelt}) \\
\text{Step 4:} & \quad A_3 = \sigma_{\text{Pers-Nr} \in B_3.\text{Pers-Nr}} (\text{Ärzte}) \\
\text{Step 5:} & \quad \text{Ergebnis} = \pi_{\text{Pers-Nr, Nachname, Vorname, Fachgebiet, Rang}} \\
& \quad \quad \quad \quad (\sigma_{\text{Pers-Nr} \in A_3.\text{Pers-Nr}} (\text{Stations\_Personal}) \bowtie A_3)
\end{align*}

\subsection*{b) SQL zu dieser Abfrage}

\textbf{Lösung:}
\begin{verbatim}
SELECT DISTINCT sp.Pers_Nr, sp.Nachname, sp.Vorname, 
       a.Fachgebiet, a.Rang
FROM Stations_Personal sp
INNER JOIN Ärzte a ON sp.Pers_Nr = a.Pers_Nr
INNER JOIN behandelt b ON a.Pers_Nr = b.Pers_Nr
INNER JOIN Patient p ON b.Pat_Nr = p.Pat_Nr
INNER JOIN Zimmer z ON p.Raum_Nr = z.Raum_Nr
WHERE z.Stat_Nr = 3;
\end{verbatim}

\vspace{2cm}

\section*{Aufgabenteil 4 - Normalisierung (9 Punkte)}

\textbf{Gegeben:}
\begin{itemize}
    \item Relationsschema [R] mit Attributen: a, b, c, d, e
    \item Funktionale Abhängigkeiten:
    \begin{itemize}
        \item $ad \to bc$
        \item $b \to ac$
        \item $d \to e$
    \end{itemize}
\end{itemize}

\subsection*{a) Schlüsselkandidaten}

\textbf{Bestimmung der Schlüsselkandidaten:}

Wir suchen minimale Mengen von Attributen, die alle anderen Attribute funktional bestimmen.

\textbf{Kandidat 1: \{a, d\}}
\begin{itemize}
    \item $ad \to bc$ (gegeben) → Bestimmt b und c \checkmark
    \item $d \to e$ → Bestimmt e \checkmark
    \item $ad$ bestimmt: a, b, c, d, e \checkmark
\end{itemize}

\textbf{Kandidat 2: \{b, d\}}
\begin{itemize}
    \item $b \to ac$ (gegeben) → Bestimmt a und c \checkmark
    \item $d \to e$ → Bestimmt e \checkmark
    \item Von a,d können wir $ad \to bc$ verwenden → Bestimmt b (redundant) und c \checkmark
    \item $bd$ bestimmt: a, b, c, d, e \checkmark
\end{itemize}

\textbf{Überprüfung auf Minimalität:}
\begin{itemize}
    \item $\{a\}$: Bestimmt nur a, nicht ausreichend
    \item $\{d\}$: Bestimmt nur d, e; nicht ausreichend
    \item $\{b\}$: Bestimmt a, b, c; nicht ausreichend (e fehlt)
    \item $\{a, d\}$ und $\{b, d\}$: Beide minimal
\end{itemize}

\textbf{Schlüsselkandidaten:} $\boxed{\{a, d\} \text{ und } \{b, d\}}$

\subsection*{b) Prüfung auf 3NF/BCNF}

\textbf{Annahme: Primärschlüssel = \{a, d\}}

\textbf{Prim-Attribute:} a, d
\textbf{Non-Prime-Attribute:} b, c, e

\textbf{3NF-Prüfung:}

Für 3NF müssen folgende Bedingungen erfüllt sein:
\begin{enumerate}
    \item Relation ist in 2NF (impliziert durch 1NF)
    \item Kein Non-Prime-Attribut hängt transitiv vom Schlüssel ab
\end{enumerate}

\textbf{Analyse der FDs:}
\begin{itemize}
    \item $ad \to bc$: Linke Seite ist Schlüssel → OK für 3NF (partielle FD)
    \item $b \to ac$: Linke Seite ist Non-Prime-Attribut
    \begin{itemize}
        \item $b$ ist nicht Kandidatenschlüssel
        \item Rechte Seite enthält a (Prime-Attribut) und c (Non-Prime-Attribut)
        \item Verletzung: Non-Prime-Attribut c hängt von Non-Prime-Attribut b ab
        \item \textbf{NICHT in 3NF!}
    \end{itemize}
    \item $d \to e$: Linke Seite ist Teil des Schlüssels
    \begin{itemize}
        \item Dies ist eine partielle Abhängigkeit des Schlüssels
        \item NICHT in 2NF! (Zu einer Non-Prime-Attribut von einem Teil des Schlüssels)
    \end{itemize}
\end{itemize}

\textbf{Ergebnis:} Die Relation ist \textbf{NICHT in 3NF} und \textbf{NICHT in 2NF}.

\subsection*{c) Zerlegung}

\textbf{Zerlegungsschritte:}

\textbf{Schritt 1: Beseitigung der Abhängigkeit von d → e}

Erstelle $R_1$ mit der FD $d \to e$:
\begin{itemize}
    \item $R_1(d, e)$ mit Schlüssel \{d\}
    \item $R_2(a, b, c, d)$ mit verbleibenden Attributen
\end{itemize}

\textbf{Schritt 2: Beseitigung der Abhängigkeit von $b \to ac$ in $R_2$}

Erstelle $R_3$ mit der FD $b \to ac$:
\begin{itemize}
    \item $R_3(b, a, c)$ mit Schlüssel \{b\}
    \item $R_4(b, d)$ mit verbleibenden Attributen
\end{itemize}

\textbf{Finale Zerlegung:}
\begin{align*}
R_1 &: (d, e) \quad \text{Schl\"ussel: } \{d\} \\
R_2 &: (b, a, c) \quad \text{Schl\"ussel: } \{b\} \\
R_3 &: (a, d) \quad \text{Schl\"ussel: } \{a, d\} \text{ (oder } \{b, d\} \text{)}
\end{align*}

\textbf{Verifizierung der 3NF:}
\begin{itemize}
    \item $R_1$: Nur $d \to e$; alle Attribute sind Schlüssel → OK
    \item $R_2$: Nur $b \to ac$; $b$ ist Schlüssel → OK
    \item $R_3$: Nur Primärattribute → OK
\end{itemize}

\subsection*{d) Verlustlosigkeit der Zerlegung}

\textbf{Zu zeigen:} $R = R_1 \bowtie R_2 \bowtie R_3$

\textbf{Nachweis durch Tupel-Verfolgung:}

Angenommen, wir haben ein Tupel $(a_0, b_0, c_0, d_0, e_0)$ in R mit:
\begin{itemize}
    \item $d_0 e_0$ in $R_1$
    \item $b_0 a_0 c_0$ in $R_2$
    \item $a_0 d_0$ in $R_3$
\end{itemize}

\textbf{Rekombination:}
\begin{align*}
R_1 \bowtie R_3 &= \{(a_0, d_0, e_0)\} \\
(R_1 \bowtie R_3) \bowtie R_2 &= \{(a_0, d_0, e_0)\} \bowtie \{(b_0, a_0, c_0)\} \\
&= \{(a_0, b_0, c_0, d_0, e_0)\}
\end{align*}

Die Rekombination ergibt genau das ursprüngliche Tupel zurück.

\textbf{Formaler Nachweis:}
Die Zerlegung ist verlustlos, wenn mindestens eine der folgenden Bedingungen erfüllt ist:
\begin{itemize}
    \item Die Zerlegung folgt einer funktionalen Abhängigkeit (Beziehung $ \to $)
    \item Die Zerlegung kann durch Chase-Algorithmus verifiziert werden
\end{itemize}

Da wir nach funktionalen Abhängigkeiten zerlegt haben ($ d \to e$ und $b \to ac$), ist die Zerlegung verlustlos.

\textbf{Ergebnis:} \textbf{Die Zerlegung ist verlustlos.} \checkmark

\end{document}
